\section{Ejemplos diversos}
\label{VI.11}

\begin{enumerate}

\item[a)] \textbf{Categorías de flechas de}~$\cal{E}$. Sea
$\cal{E}$ una categoría. Denotemos por $\mathbf{\Delta^1}$ la categoría
asociada al conjunto totalmente ordenado de dos elementos $[0,1]$;
\marginpar{183}
tiene, pues, dos objetos, 0 y 1, y, además de los dos morfismos identidad,
una flecha $(0,1)$ con fuente 0 y destino 1. Sea
$$
\CatFl(\cal{E})=\SheafHom(\mathbf{\Delta^1},\cal{E});
$$
se llama la \emph{categoría de flechas de}~$\cal{E}$. El objeto 1
de~$\mathbf{\Delta^1}$ define un funtor canónico, llamado
\emph{funtor codominio},
$$
\CatFl(\cal{E})\to \cal{E}
$$
(el funtor definido por el objeto 0 de~$\mathbf{\Delta^1}$ se llama
\emph{funtor dominio}). Para todo objeto $S$ de~$\cal{E}$, la categoría
fibra $\CatFl(\cal{E})_S$ es canónicamente isomorfa a la categoría
$\cal{E}_{/S}$ de los objetos de~$\cal{E}$ sobre~$S$.

Consideremos un morfismo $f\colon T\to S$ de $\cal{E}$. Le corresponde
entonces un funtor canónico
$$
f_*\colon \cal{E}_{/T}=\cal{F}_T\to \cal{E}_{/S}=\cal{F}_S
$$
y un isomorfismo natural
$$
\Hom_S(f_*(\eta),\xi)\isomto\Hom_f(\eta,\xi)
$$
que hace de~$f_*$ un funtor de imagen directa para $f$ en~$\cal{F}$.
Además, en este caso se tiene
$$
(\id_S)_*=\id_{\cal{F}_S},\quad (fg)_*=f_*g_*,\quad
\ifthenelse{\boolean{orig}}
{c^{f,g}=\id_{(fg)},}
{c^{f,g}=\id_{(fg)_*},}
$$
es decir, $\cal{F}$ está dotada de una coescisión sobre~$\cal{E}$. Con
mayor razón, $\cal{F}$ es cofibrada sobre~$\cal{E}$. Observemos ahora que
el conjunto de morfismos de $\cal{F}$ está en correspondencia biunívoca
con el conjunto de diagramas cuadrados conmutativos de~$\cal{E}$:
$$
\xymatrix{
X\ar[d]_-u&\ar[l]_-{f'}Y\ar[d]^-v\\
S&\ar[l]_-{f}T
}
$$
Por
\marginpar{184}
definición, el morfismo en cuestión es cartesiano si el cuadrado es
cartesiano en~$\cal{E}$, es decir, si hace de~$Y$ un producto fibrado de
$X$ y $T$ sobre~$S$. Por tanto, el funtor de imagen inversa
\ifthenelse{\boolean{orig}}
{$f$}
{$f^*$}
existe si y solo si para todo objeto $X$ sobre~$S$ existe el producto
fibrado $X\times_ST$. Se deduce de~\Ref{VI.10.1} que, si en~$\cal{E}$
existe siempre el producto de dos objetos sobre un tercero, es decir, si
$\cal{F}$ es prefibrada sobre~$\cal{E}$, entonces $\cal{F}$ es incluso
fibrada sobre~$\cal{E}$.

\medskip
\item[b)] \textbf{Categoría de prehaces o haces sobre espacios variables}

Sea $\cal{E}=\mathbf{Top}$ la categoría de los espacios topológicos. Si
$T$ es un espacio topológico, denotaremos por $\cal{U}(T)$ la categoría de
los abiertos de~$T$, cuyos morfismos son las aplicaciones de inclusión. Si
$\cal{C}$ es una categoría, un funtor
$\cal{U}(T)^\circ\to \cal{C}$ se llama un \emph{prehaz} sobre~$T$ con
valores en~$\cal{C}$, y un \emph{haz} si satisface una condición de
exactitud por la izquierda que no repetiremos aquí. La \emph{categoría
$\cal{P}(T)$ de prehaces sobre~$T$ con valores en~$\cal{C}$} es, por
definición, la categoría $\SheafHom(\cal{U}(T)^\circ,\cal{C})$, y la
categoría $\cal{F}(T)$ de haces sobre~$T$ con valores en~$\cal{C}$ es la
subcategoría plena cuyos objetos son los objetos de
$\SheafHom(\cal{U}(T)^\circ,\cal{C})$ que son haces. Si
$f\colon T\to S$ es un morfismo de~$\cal{E}$, es decir, una aplicación
continua de espacios topológicos, la aplicación creciente
$U\mto f^{-1}(U)$ le hace corresponder un funtor
$\cal{U}(S)\to\cal{U}(T)$ y, por tanto, un funtor
$$
f_*\colon\SheafHom(\cal{U}(T)^\circ,\cal{C})\to
\SheafHom(\cal{U}(S)^\circ,\cal{C}),
$$
llamado \emph{funtor de imagen directa de prehaces por}~$f$. Se ve de
inmediato que la imagen directa de un haz es un haz; por tanto, el funtor
\ifthenelse{\boolean{orig}}
{$f_*\cal{P}(T)\to\cal{P}(S)$}
{$f_*\colon \cal{P}(T)\to\cal{P}(S)$}
induce un funtor, denotado también
\ifthenelse{\boolean{orig}}
{$f_*\colon\cal{P}(T)\to\cal{P}(S)$.}
{$f_*\colon\cal{F}(T)\to\cal{F}(S)$.}
Además, se comprueba trivialmente, por la asociatividad de la composición
de funtores, que para una segunda aplicación continua $g\colon U\to T$ se
tiene la identidad
$$
\ifthenelse{\boolean{orig}}
{(gf)_* =g_*f_*,}
{(fg)_*=f_*g_*,}
\qquad
\text{y asimismo}\quad (\id_S)_*=\id_{\cal{P}(S)}.
$$
De esta forma se ha obtenido un funtor
$$
S\mto \cal{P}(S)
$$
\resp $$
S\mto \cal{F}(S)
$$
de
\marginpar{185}
$\cal{E}$ en~$\Cat$. En realidad, nos interesa el funtor correspondiente
$$
S\mto \cal{P}(S)^\circ,\qquad\resp\qquad S\mto
\cal{F}(S)^\circ.
$$
Define una categoría cofibrada, e incluso coescindida, sobre la categoría
de los espacios topológicos, que se llama la \emph{categoría cofibrada de
los prehaces} (\resp \emph{haces}) \emph{con valores en}~$\cal{C}$
(se sobreentiende: sobre espacios variables). Explicitando la construcción
del \No~\Ref{VI.8}, se ve que un morfismo de un prehaz $B$ sobre~$T$ en un
prehaz $A$ sobre~$S$ es un par $(f,u)$ formado por una aplicación continua
de~$T$ en~$S$ y un morfismo $u\colon A\to f_*(B)$ en la categoría
$\cal{P}(S)$. Esta descripción vale también para los morfismos de haces,
pues $\cal{F}$ es una subcategoría plena de~$\cal{P}$.

En los casos más importantes, la categoría $\cal{P}$ y la categoría
$\cal{F}$ sobre~$\cal{E}$ son también categorías fibradas, es decir, para
toda aplicación continua los funtores de imagen directa
$\cal{P}(T)\to\cal{P}(S)$ y $\cal{F}(T)\to\cal{F}(S)$ tienen un funtor
adjunto, que se denota entonces por $f^*$ y se llama funtor de imagen
inversa de prehaces \resp funtor de imagen inversa de haces por la
aplicación continua~$f$. Este funtor existe, por ejemplo, si
$\cal{C}=\Ens$. Se puede demostrar que el funtor
$f^*\colon\cal{P}(S)\to\cal{P}(T)$ existe siempre que en~$\cal{C}$
existan los límites inductivos relativos a diagramas del universo
considerado. La cuestión es menos fácil para $\cal{F}$; en efecto, incluso
en el caso $\cal{C}=\Ens$, la imagen inversa de un prehaz que es un haz no
es, en general, un haz. En otras palabras, el funtor de imagen inversa de
haces no es isomorfo al funtor inducido por el funtor de imagen inversa de
prehaces, pese a la notación común $f^*$. Así, $\cal{F}$ es una
subcategoría cofibrada de~$\cal{P}$, pero no una subcategoría fibrada, es
decir, \emph{el funtor de inclusión $\cal{F}\to\cal{P}$ no es fibrante}.

La categoría cofibrada $\cal{P}$ puede deducirse de una categoría
cofibrada, o más bien fibrada, más general, obtenida como sigue. Para toda
categoría $\cal{U}$ del universo fijado, pongamos
$$
\cal{P}(\cal{U})=\SheafHom(\cal{U},\cal{C})
$$
y
\marginpar{186}
observemos que $\cal{U}\mto\cal{P}(\cal{U})$ es de manera natural un
funtor contravariante en~$\cal{U}$, de la categoría $\Cat$ en~$\Cat$.
Define, pues, una categoría escindida sobre~$\cal{E}=\Cat$, que denotaremos
por $\Cat_{/\!/\cal{C}}$. Los objetos de esta categoría son los pares
$(\cal{U},p)$ formados por una categoría $\cal{U}$ y un funtor
$p\colon\cal{U}\to\cal{C}$, y un morfismo de $(\cal{U},p)$ en
$(\cal{V},q)$ es esencialmente un par $(f,u)$, donde $f$ es un funtor
$\cal{U}\to\cal{V}$ y $u$ una transformación natural de funtores
$u\colon p\to qf$. Dejamos al lector la tarea de explicitar la composición
de los morfismos en $\Cat_{/\!/\cal{C}}$. El funtor de proyección
$$
\cal{F}=\Cat_{/\!/\cal{C}}\to\cal{E}=\Cat
$$
asocia al par $(\cal{U},p)$ el objeto $\cal{U}$; la categoría fibra en
$\cal{U}$ es la categoría $\SheafHom(\cal{U},\cal{C})$, salvo isomorfismo.
Cuando existen en~$\cal{C}$ los límites inductivos, se demuestra fácilmente
que la categoría fibrada $\Cat_{/\!/\cal{C}}$ sobre~$\Cat$ es también
cofibrada sobre~$\Cat$, es decir, se puede definir la noción de
\emph{imagen directa de un funtor} $p\colon\cal{U}\to\cal{C}$ por un
funtor $f\colon\cal{U}\to\cal{V}$. La categoría de prehaces se deduce de
la categoría fibrada precedente mediante el cambio de base
$$
\mathbf{Top}^\circ\to\Cat
$$
(el funtor $S\mto\cal{U}(S)$ definido más arriba), lo que da una categoría
fibrada sobre~$\mathbf{Top}^\circ$; al pasar a la categoría opuesta se
obtiene la categoría cofibrada $\cal{P}$ de los prehaces sobre
$\mathbf{Top}$. La noción de imagen inversa de un funtor corresponde a la
de imagen directa de un prehaz, y la noción de imagen directa de un funtor,
a la de imagen inversa de un prehaz.

\medskip
\item[c)] \textbf{Objetos con operadores sobre un objeto con operadores}

Sea $\cal{F}$ una categoría sobre~$\cal{E}$, y sea $S$ un objeto de
$\cal{E}$ sobre el que actúa un grupo $G$, por la izquierda para fijar
ideas. Este objeto con operadores puede interpretarse como correspondiente
a un funtor $\lambda\colon\cal{E}'\to\cal{E}$ desde la categoría (con un
único objeto y con $G$ como grupo de endomorfismos) $\cal{E}'$, definida
por~$G$, a la categoría $\cal{E}$. Por cambio de base define,
por tanto, una categoría $\cal{F}'$ sobre~$\cal{E}'$, que es fibrada
\resp cofibrada cuando $\cal{F}$ lo es sobre~$\cal{E}$. Una
\marginpar{187}
\ifthenelse{\boolean{orig}}
{sección de~$\cal{E}'$ sobre~$\cal{F}'$}
{sección de~$\cal{F}'$ sobre~$\cal{E}'$}
(necesariamente cartesiana, pues $\cal{E}'$ es un grupoide y todo
isomorfismo de~$\cal{F}'$ es cartesiano en virtud de~\Ref{VI.6.12}) puede
interpretarse también como un $\cal{E}$-funtor
$\cal{E}'\to\cal{F}$ sobre~$\lambda$, o como un objeto con operadores
$\xi$ en~$\cal{F}$ «sobre» el objeto con operadores~$S$.

\medskip
\item[d)] \textbf{Pares de funtores adjuntos cuasi-inversos;
autodualidades}

Cuando la categoría base $\cal{E}$ se reduce a dos objetos $a,b$ y, además
de las flechas identidad, a dos isomorfismos $f\colon a\to b$ y
$g\colon b\to a$ inversos entre sí, es decir, cuando $\cal{E}$ es un
grupoide conexo rígido con dos objetos, una categoría con clivaje
normalizado sobre~$\cal{E}$ es esencialmente lo mismo que el sistema
formado por dos categorías $\cal{F}_a$ y $\cal{F}_b$ y un \emph{par de
funtores adjuntos} $G\colon\cal{F}_a\to\cal{F}_b$ y
$F\colon\cal{F}_b\to\cal{F}_a$ que sean equivalencias de categorías y,
por tanto, cuasi-inversos entre sí. Se toman como $\cal{F}_a$ y
$\cal{F}_b$ las categorías fibra de~$\cal{F}$, como $F$ y $G$ los funtores
$f^*$ y $g^*$, y los dos isomorfismos
$$
u\colon FG\isomto\id_{\cal{F}_a}\qquad v\colon
GF\isomto\id_{\cal{F}_b}
$$
son $c_{g,f}$ y $c_{f,g}$. Las dos condiciones usuales de compatibilidad
entre $u$ y $v$ no son sino la condición \Ref{VI.7.4}~$B)$ para los
compuestos $fgf$ y $gfg$. Es fácil demostrar que estas condiciones bastan
para implicar que se tiene efectivamente un pseudofuntor
$\cal{E}^\circ\to\Cat$.

Un caso interesante es aquel en que
$$
\cal{F}_b=\cal{F}_a^\circ,\quad G=F^\circ,\quad v=u^\circ.
$$
Se llama \emph{autodualidad}
\index{autodualidad|hyperpage}%
en una categoría $\cal{C}$ al dato de un funtor
$D\colon\cal{C}\to\cal{C}^\circ$ y de un isomorfismo
\ifthenelse{\boolean{orig}}
{$u\colon DD^\circ\isomto\id_{\cal{C}}$,}
{$u\colon D^\circ D\isomto\id_{\cal{C}}$,}
tales que $u$ y el isomorfismo
\ifthenelse{\boolean{orig}}
{$u^\circ\colon D^\circ D\isomto\id_{\cal{C}^\circ}$}
{$u^\circ\colon DD^\circ\isomto\id_{\cal{C}^\circ}$}
hagan de $(D,D^\circ)$ un par de funtores adjuntos, necesariamente
cuasi-inversos entre sí. Esta condición se escribe
$$
D(u(x))=u(D(x))\qquad \text{para todo $x\in\Ob(\cal{C})$.}
$$

\item[e)]
\marginpar{188}
\textbf{Categorías sobre una categoría discreta $\cal{E}$.} Se dice que
$\cal{E}$ es una \emph{categoría discreta} si toda flecha de ella es una
flecha identidad, de modo que $\cal{E}$ queda definida, salvo isomorfismo
único, por el conjunto $I=\Ob(\cal{E})$. El dato de una categoría
$\cal{F}$ sobre~$\cal{E}$ equivale, pues, salvo isomorfismo único, al dato
de una familia de categorías $\cal{F}_i$ ($i\in I$), las categorías fibra.
Toda categoría $\cal{F}$ sobre~$\cal{E}$ es fibrada, todo $\cal{E}$-funtor
$\cal{F}\to\cal{G}$ es cartesiano y se tiene un isomorfismo canónico
$$
\SheafHom_{\cal{E}/-}(\cal{F},\cal{G}) \isomto
\underset{i}\prod\SheafHom(\cal{F}_i,\cal{G}_i).
$$
En particular, se obtiene
$$
\bf{\Gamma}(\cal{F}/\cal{E})=\varprojLim\cal{F}/\cal{E}\isomto
\ifthenelse{\boolean{orig}}
{\underset{i}\prod \cal{E}_i.}
{\underset{i}\prod \cal{F}_i.}
$$

\item[f)] \textbf{Supongamos que $\cal{E}$ tiene exactamente dos objetos
$S$ y $T$,} y, además de los \textbf{morfismos identidad, un morfismo}
$f\colon T\to S$. Entonces una categoría $\cal{F}$ sobre~$\cal{E}$ queda
definida, salvo $\cal{E}$-isomorfismo único, por el dato de dos categorías
$\cal{F}_S$ y $\cal{F}_T$ y de un bifuntor $H(\eta,\xi)$ sobre
$\cal{F}_T^\circ\times\cal{F}_S$ con valores en~$\Ens$. En efecto, si
$\cal{F}$ es una categoría sobre~$\cal{E}$, se le asocian las dos
categorías fibra $\cal{F}_S$ y $\cal{F}_T$ y el bifuntor
$H(\eta,\xi)=\Hom_f(\eta,\xi)$. Dejamos al lector la tarea de explicitar la
construcción en sentido inverso. Para que la categoría considerada sea
fibrada, o prefibrada, lo que equivale a lo mismo, es necesario y
suficiente que el funtor $H$ sea representable respecto del argumento
$\xi$. Para que sea cofibrada, es necesario y suficiente que $H$ sea
representable respecto del argumento~$\eta$.

\item[g)] Sea $\cal{F}=\cal{C}\times\cal{E}$, considerada como categoría
sobre~$\cal{E}$ mediante~$\pr_2$. Entonces $\cal{F}$ es fibrada y
cofibrada sobre~$\cal{E}$, e incluso está dotada de una escisión y una
coescisión canónicas, correspondientes al funtor constante sobre
$\cal{E}$, \resp sobre~$\cal{E}^\circ$, con valores en~$\Cat$ y valor
$\cal{C}$. Se tiene
$$
\bf{\Gamma}(\cal{F}/\cal{E})\simeq\SheafHom(\cal{E},\cal{C})
$$
y
\marginpar{189}
$\varprojLim\cal{F}/\cal{E}$ corresponde a la subcategoría plena formada
por los funtores $F\colon\cal{E\to\cal{C}}$ que transforman morfismos
cualesquiera en isomorfismos.
\end{enumerate}
